\documentclass[10pt,a4paper]{article}

\usepackage[utf8]{inputenc}
\usepackage[T1]{fontenc}
\usepackage{lmodern}
\input{glyphtounicode}
\pdfgentounicode=1
\usepackage[margin=0.52in]{geometry}
\usepackage[hidelinks]{hyperref}
\usepackage{enumitem}

\pagestyle{empty}
\setlength{\parindent}{0pt}
\setlength{\parskip}{0pt}
\setlist[itemize]{leftmargin=1.1em,itemsep=1pt,topsep=1pt,parsep=0pt}
\urlstyle{same}

\newcommand{\cvsection}[1]{%
    \vspace{0.38em}
    {\large\bfseries #1}\par
    \vspace{0.08em}
    \hrule
    \vspace{0.18em}
}

\newcommand{\entry}[4]{%
    \textbf{#1} \hfill \textbf{#2}\par
    \textit{#3}\par
    #4
    \vspace{0.14em}
}

\begin{document}

{\LARGE\bfseries ANDERSON H. R. FERREIRA}\par
\vspace{0.18em}
{\normalsize Dados, Machine Learning e Analytics | Doutorando UNICAMP}\par
\vspace{0.28em}

Salto, SP | +55 11 95850-9690 | \href{mailto:a058899@dac.unicamp.br}{a058899@dac.unicamp.br}\par
\href{https://anderson-ferreira-83.github.io/}{Portfolio} | \href{https://www.linkedin.com/in/anderson-ferreira-1a473138/?locale=pt}{LinkedIn} | \href{https://github.com/anderson-ferreira-83}{GitHub}\par
Defesa prevista: dez/2026 | Disponibilidade hibrida em Sao Paulo entre jun/2026 e set/2026

\cvsection{Resumo}
Doutorando em Engenharia Mecanica na UNICAMP, com atuacao aplicada em dados, machine learning, cloud e estatistica. Experiencia em modelagem preditiva, validacao de modelos, feature engineering e implementacao com Python, SQL, FastAPI, Oracle e AWS. Repertorio construido em projetos de IoT e saude, com tecnicas transferiveis para analytics, risco, fraude, churn e apoio a decisao.

\cvsection{Competencias Tecnicas}
Python, SQL, Java, JavaScript, FastAPI, Oracle SQL, AWS, Docker, scikit-learn, Random Forest, SMOTE, SHAP, ETL, EDA, Feature Engineering, Cross-validation, Holdout, Threshold Analysis, Estatistica Aplicada.

\cvsection{Projetos de Dados e ML}
\entry{Predicao de Risco com SMOTE}{Projeto aplicado}{Saude | orientador e desenvolvedor}{
\begin{itemize}
    \item Atuei em projeto de predicao de risco de hipertensao com 4.240 amostras e 12 variaveis clinicas.
    \item Implementei split estratificado, SMOTE apenas no treino e validacao cruzada 5-fold, evitando data leakage em problema desbalanceado.
    \item Estruturei avaliacao com foco em Recall, F2-score, AUC-ROC, thresholds clinicos e interpretabilidade com SHAP.
    \item Abordagem transferivel para fraude, churn, credit scoring e outros cenarios com classe rara.
\end{itemize}
}

\entry{Pipeline IoT End-to-End}{Projeto aplicado}{ESP32 + FastAPI + Oracle XE}{
\begin{itemize}
    \item Desenvolvi pipeline de coleta, tratamento e classificacao para 7 estados operacionais com ESP32, MPU6050, FastAPI e Oracle XE.
    \item Expandi o espaco de atributos para 104 features candidatas e selecionei 16 finais; Random Forest com 200 arvores atingiu 97,34\% +- 0,63\% em validacao cruzada e 94,24\% em holdout.
    \item Projeto com mais de 210 mil registros em runs documentadas e inferencia no browser sem servidor de ML dedicado.
\end{itemize}
}

\cvsection{Experiencia Profissional}
\entry{Professor Adjunto}{fev/2025--Atual}{Cruzeiro do Sul Educacional}{
\begin{itemize}
    \item Docencia em Programacao Orientada a Objetos, Estrutura de Dados, Programacao Web e Fundamentos de Bancos de Dados.
    \item Desenvolvimento de atividades praticas com Java, Python e SQL, conectando algoritmos, modelagem de dados e aplicacoes reais.
    \item Acompanhamento de projetos aplicados com integracao entre software, dados e bancos de dados.
\end{itemize}
}

\entry{Pesquisador de Doutorado}{2018--Atual}{UNICAMP - Engenharia Mecanica}{
\begin{itemize}
    \item Pesquisa aplicada em processamento de sinais, modelagem estatistica, sensores e classificacao de estados operacionais.
    \item Publicacao aceita no Mechanical Systems and Signal Processing (MSSP), 2026, comparando 5 geometrias com abordagem orientada a dados.
\end{itemize}
}

\entry{Professor Universitario}{2014--2021}{CEUNSP}{
\begin{itemize}
    \item Atuei em disciplinas de Probabilidade e Estatistica Basica, Fisicas Experimentais, Termodinamica Aplicada e Transferencia de Calor.
    \item Desenvolvi atividades ligadas a regressao, testes de hipotese, inferencia estatistica e simulacao computacional.
\end{itemize}
}

\cvsection{Formacao}
\textbf{Doutorado em Engenharia Mecanica} - UNICAMP \hfill 2018--2026\par
\textbf{Mestrado em Engenharia Mecanica} - UNICAMP \hfill 2011--2013\par
\textbf{Licenciatura em Fisica} - UNICAMP \hfill 2006--2010

\cvsection{Certificacoes e Publicacao}
AWS Certified AI Practitioner | Applied Machine Learning in Python | IBM Data Engineering Professional Certificate\par
Mechanical Systems and Signal Processing (MSSP), 2026 | Impact Factor 8.9

\end{document}
